\documentclass{beamer}
\usetheme{Madrid}
\usecolortheme{default}
\usepackage{fontspec}
\usepackage{polyglossia}
\setdefaultlanguage{english}
\usepackage{tikz}
\usepackage{hyperref}
\usepackage{graphicx}
\usepackage{booktabs}

% Custom colors
\definecolor{RurDynamixGreen}{RGB}{76, 175, 80}
\definecolor{RurDynamixBlue}{RGB}{33, 150, 243}
\definecolor{DarkGray}{RGB}{64, 64, 64}

\setbeamercolor{primary}{fg=white,bg=RurDynamixGreen}
\setbeamercolor{secondary}{fg=white,bg=RurDynamixBlue}
\setbeamertemplate{frametitle}{
  \begin{beamercolorbox}[rounded=true,shadow=true,leftskip=1cm,rightskip=1cm,wd=\paperwidth,ht=0.8cm,dp=0.4cm]{frametitle}
    \usebeamerfont{frametitle}\insertframetitle
  \end{beamercolorbox}
}

% Title Page Setup
\title{\tiny{Rural Territories in Transition: Morocco and the Global South}}
\author{Said Moukrim}
\institute{UM5-FSR Morocco}
\date{November 28, 2026}

\begin{document}

\begin{frame}

\begin{columns}[T]
  \column{0.5\textwidth}
  \includegraphics[width=\textwidth]{images/rur_dyn.png}

  \column{0.25\textwidth}
  \includegraphics[width=\textwidth]{images/FSR logo.png}
\end{columns}

\vspace{1cm}

\begin{center}
\Large {Hydraulic Innovation and Rural Transitions:}\\
\vspace{0.8cm}
\large{Artificial Basins Dynamics in Bab Berred}\\
Central Rif (Morocco)\\
\vspace{1.5cm}
\Large{Said Moukrim}
\end{center}

\end{frame}

% ==================== INTRODUCTION ====================

% Slide 2: Problem
\begin{frame}
\frametitle{Problem: Water Scarcity in Mountains}

\begin{itemize}
  \item Precipitation is \textbf{declining} in Bab Berred
  \item \textbf{2023}: Record low = 300 mm (very dry)
  \item Mountains: Water is essential for agriculture
  \item Farmers need solutions
\end{itemize}

\vspace{0.5cm}

\textbf{Farmers' response:}
\begin{itemize}
  \item Build artificial water basins
  \item Store water from rainy season
  \item Use water in dry season
\end{itemize}

\end{frame}

% Slide 3: Why Cannabis? (Market, not law)
\begin{frame}
\frametitle{Why Cannabis? Market Demand, Not Law}

\textbf{Important correction:}
\begin{itemize}
  \item \textbf{NOT}: Legalization in 2021 caused the expansion
  \item \textbf{YES}: International market demand caused the expansion
  \item Law 13-21 only allowed official channels
\end{itemize}

\vspace{0.5cm}

\textbf{Why farmers choose hybrid cannabis:}
\begin{itemize}
  \item Old variety (kif): 15 kg/hectare
  \item New hybrids: 70 kg/hectare
  \item \textbf{4-5× more productive!}
  \item Farmers earn more money
\end{itemize}

\vspace{0.5cm}
\textbf{The cost}: Hybrids need much more water $\Rightarrow$ basins

\end{frame}

% Slide 4: Study Area
\begin{frame}
\frametitle{Study Area: Bab Berred}

\textbf{Where is it?}
\begin{itemize}
  \item Chefchaouen Province, Morocco
  \item Mountain region (Rif Mountains)
  \item Population: 24,742 people
\end{itemize}

\vspace{0.3cm}

\textbf{Terrain:}
\begin{itemize}
  \item Very steep: 80\% average slope
  \item High altitude: 650-2,100 meters
  \item Total area: 172.24 km²
\end{itemize}

\vspace{0.3cm}

\begin{center}
  \vspace{1.5cm}
  \includegraphics[width=0.7\textwidth]{images/Figure1_Location.png}
  \small{\textit{Figure 1: Location of Bab Berred}}
\end{center}

\end{frame}

% ==================== METHODS ====================

% Slide 5: Methods
\begin{frame}
\frametitle{How We Did This Study}

\textbf{Data sources:}
\begin{itemize}
  \item \textbf{Satellite images}: Landsat (1998-2025) = see vegetation
  \item \textbf{Google Earth Pro}: See basins (manual counting)
  \item \textbf{Precipitation data}: From satellites (1985-2025)
  \item \textbf{Land-use maps}: Classify vegetation types
\end{itemize}

\vspace{0.3cm}

\textbf{What we measured:}
\begin{itemize}
  \item How many basins? When built?
  \item Where are they located?
  \item How big are they?
  \item How is vegetation changing? (NDVI)
  \item How is precipitation changing?
\end{itemize}

\vspace{0.3cm}

\textbf{Tools:}
\begin{itemize}
  \item QGIS (mapping software)
  \item Google Earth Engine (satellite analysis)
  \item RStudio (statistics)
\end{itemize}

\end{frame}

% ==================== RESULTS ====================

% Slide 6: Basin Numbers
\begin{frame}
\frametitle{Result 1: Basin Explosion}

\textbf{Basin count:}
\begin{itemize}
  \item 2010: $\approx$ 28 basins
  \item 2025: 2,044 basins
  \item Growth: 73$\times$ more basins in 15 years!
\end{itemize}

\vspace{0.3cm}

\textbf{Storage capacity:}
\begin{itemize}
  \item 2010: 50,000 m³ total water stored
  \item 2025: 4.3 million m³ total
  \item Increase: 85$\times$ more water!
\end{itemize}

\vspace{0.3cm}

\begin{center}
  \vspace{1.2cm}
  \includegraphics[width=0.75\textwidth]{images/Figure2_BasinGrowth.png}
  \small{\textit{Figure 2: Number of basins per year (2010-2025)}}
\end{center}

\end{frame}

% Slide 7: Where Are Basins?
\begin{frame}
\frametitle{Result 2: Where Are Basins Located?}

\textbf{Altitude preference:}
\begin{itemize}
  \item Concentration: 800-1,500 m
  \item Not too high, not too low
\end{itemize}

\vspace{0.2cm}

\textbf{Slope preference:}
\begin{itemize}
  \item 71\% on slopes 5-15\%
  \item Avoid very steep slopes (too expensive)
\end{itemize}

\vspace{0.2cm}

\textbf{Land-use:}
\begin{itemize}
  \item 78\% in cannabis fields
  \item 18\% in grasslands
  \item 3\% in natural forest
\end{itemize}

\vspace{0.3cm}

\begin{center}
  \vspace{0.8cm}
  \includegraphics[width=0.75\textwidth]{images/Figure3_SpatialDistribution.png}
  \small{\textit{Figure 3: Where basins are located (maps)}}
\end{center}

\end{frame}

% Slide 8: Vegetation Improvement?
\begin{frame}
\frametitle{Result 3: Vegetation Gets Greener (NDVI)}

\textbf{What is NDVI?}
\begin{itemize}
  \item Satellite measurement of vegetation (greenness)
  \item 0 = bare soil, 1 = dense forest
  \item Higher = more green
\end{itemize}

\vspace{0.3cm}

\textbf{What we see:}
\begin{itemize}
  \item 1998-2012: Slight increase
  \item 2012-2025: Big increase in greenness
  \item Peak greenness: May-August (cannabis growing season)
\end{itemize}

\vspace{0.3cm}

\begin{center}
  \vspace{0.8cm}
  \includegraphics[width=0.75\textwidth]{images/Figure4_NDVI.png}
  \small{\textit{Figure 4: NDVI over time (green = higher values after 2012)}}
\end{center}

\end{frame}

% Slide 9: But Precipitation Declines!
\begin{frame}
\frametitle{Result 4: Precipitation DECLINES}

\textbf{The problem:}
\begin{itemize}
  \item Overall trend: Precipitation is declining
  \item Statistical test: tau = -0.1376 (significant decline)
  \item Year-to-year: Big ups and downs
  \item 2023: Record low ($\approx$ 300 mm)
\end{itemize}

\vspace{0.3cm}

\begin{center}
  \vspace{1.2cm}
  \includegraphics[width=0.75\textwidth]{images/Figure5_Precipitation.png}
  \small{\textit{Figure 5: Annual precipitation (declining trend visible)}}
\end{center}

\vspace{0.3cm}

\small{\textit{The paradox: Vegetation greens BUT rain decreases!}}

\end{frame}

% ==================== DISCUSSION ====================

% Slide 10: The Paradox Explained
\begin{frame}
\frametitle{The Paradox: Why NDVI Up When Rain Down?}

\textbf{Answer: Irrigation!}
\begin{itemize}
  \item Basins store water from rainy seasons
  \item Water used for irrigation in dry season
  \item Cannabis gets watered $\Rightarrow$ green vegetation
  \item Satellites see green $\Rightarrow$ high NDVI
\end{itemize}

\vspace{0.3cm}

\textbf{But is this good?}
\begin{itemize}
  \item NDVI high = lots of green vegetation
  \item BUT: Only cannabis (monoculture), not forest
  \item Natural areas (forest, grassland) = \textbf{decreased}
  \item High NDVI masks land conversion!
\end{itemize}

\vspace{0.3cm}

\textbf{Bottom line:}
\begin{itemize}
  \item More vegetation ≠ healthier ecosystem
  \item Just more farming, not more nature
\end{itemize}

\end{frame}

% Slide 11: The Real Cost
\begin{frame}
\frametitle{The Cost: Aquifer Depletion}

\textbf{Critical problem:}
\begin{itemize}
  \item Farmers extract water faster than rain refills aquifer
  \item Groundwater level = \textbf{declining}
  \item Long-term: Water will run out
\end{itemize}

\vspace{0.3cm}

\textbf{Other impacts:}
\begin{itemize}
  \item Natural habitat fragmented (by basins)
  \item Forest converted to cannabis fields
  \item Endemic species at risk
\end{itemize}

\vspace{0.3cm}

\textbf{The paradox of adaptation:}
\begin{itemize}
  \item Short-term: Basins solve water problem
  \item Long-term: Basins create new problem (aquifer depletion)
  \item Rational for farmer today $\neq$ sustainable for region tomorrow
\end{itemize}

\end{frame}

% ==================== CONCLUSION ====================

% Slide 12: Conclusion
\begin{frame}
\frametitle{Conclusion}

\textbf{What we found:}
\begin{itemize}
  \item Basins increased 73× in 15 years
  \item Vegetation looks greener (but it's just crops)
  \item Rain is declining (long-term trend)
  \item Aquifers are depleting
\end{itemize}

\vspace{0.3cm}

\textbf{Why it matters:}
\begin{itemize}
  \item Same problem in many Mediterranean countries
  \item Same problem in semi-arid regions worldwide
  \item Need better water management
  \item Infrastructure alone is not enough
\end{itemize}

\vspace{0.3cm}

\textbf{What needs to happen:}
\begin{itemize}
  \item Sustainable water limits
  \item Protect remaining natural areas
  \item Diversify farming (not just cannabis)
  \item Work together (not individual farmers separately)
\end{itemize}

\end{frame}

% Final Slide
\begin{frame}[plain]
\begin{center}
  \vfill
  {\Huge \textbf{Thank you}}
  \\[1cm]
  {\Large Questions?}
  \\[2cm]
  \small{
    Said Moukrim, Ousmane Samba MAÏGA, et al. \\
    Mohammed-V University in Rabat \\
    Faculty of Sciences
  }
  \vfill
\end{center}
\end{frame}

\end{document}
